\documentclass[11pt]{article}

\usepackage[a4paper,margin=1in]{geometry}
\usepackage[T1]{fontenc}
\usepackage{lmodern}
\usepackage{microtype}
\usepackage{amsmath,amssymb,amsthm,mathtools}
\usepackage{booktabs,tabularx}
\usepackage{enumitem}
\usepackage[numbers,sort&compress]{natbib}
\usepackage[colorlinks=true,linkcolor=blue,citecolor=blue,urlcolor=blue]{hyperref}
\usepackage[nameinlink,noabbrev]{cleveref}

\setlength{\emergencystretch}{3em}

\newtheorem{theorem}{Theorem}[section]
\newtheorem{proposition}[theorem]{Proposition}
\newtheorem{corollary}[theorem]{Corollary}
\newtheorem{lemma}[theorem]{Lemma}
\theoremstyle{definition}
\newtheorem{definition}[theorem]{Definition}
\theoremstyle{remark}
\newtheorem{remark}[theorem]{Remark}

\newcommand{\one}{\mathbf 1}
\newcommand{\Lzero}{\mathrm{L0}}
\newcommand{\Lone}{\mathrm{L1}}
\newcommand{\Ltwo}{\mathrm{L2}}

\title{\textbf{Squarefree Masks in the Shift-Two Frame}\\
\large Exact CRT Fibers, Primitive Truncation Tails, and the
Growing-Modulus Ledger}
\author{Liang Wang}
\date{July 2026}

\begin{document}
\maketitle

\begin{abstract}
TPC-127 writes a primitive determinant-two M\"obius pair as a
shift-two Liouville product on a growing progression, multiplied by
the quotient-squarefree mask
\[
 \mu^2(D(z))\mu^2(V(z)).
\]
We resolve this mask without appealing to a density heuristic.  For
\[
 Q_R(m)=\sum_{\substack{k\le R\\k^2\mid m}}\mu(k)
\]
the truncated product has the exact CRT expansion
\[
 Q_R(D(z))Q_R(V(z))
 =
 \sum_{\substack{k,\ell\le R\\
 (k,s)=(\ell,a)=(k,\ell)=1}}
 \mu(k)\mu(\ell)
 \one_{z\equiv r_{k,\ell}\ ({\rm mod}\ k^2\ell^2)}.
\]
In the TPC-127 variable \(n=sV=aD+2\), the same fiber is the unique
solution of
\[
 n\equiv2\pmod{ak^2},\qquad n\equiv0\pmod{s\ell^2},
\]
and its modulus is \(ask^2\ell^2\le asR^4\).  If the affine values
are at most \(Y\) on an interval of length \(N\), the total
unweighted truncation error is
\[
 \ll \tau_*(Y)\left(\frac NR+\sqrt Y\right),
\qquad
 \tau_*(Y)=\max_{m\le Y}\tau(m).
\]
We also prove the component census
\(\ll N+R^2\), which replaces a wasteful \(R^2N\) triangle bound.
These are unconditional finite and elementary estimates.  They
isolate, but do not prove, the required uniform cancellation on the
resulting progression fibers.  Thus the paper is \(\Lzero/\Lone\),
not an \(\Ltwo\) parity result.
\end{abstract}

\section{Primitive determinant-two setting}

Let
\begin{equation}
 D(z)=d+sz,\qquad V(z)=u+az,\qquad su-ad=2
\label{eq:forms}
\end{equation}
be positive on an integer interval \(I\) of length \(N\).  We use
the primitive relations supplied by the literal TPC-93 export,
\begin{equation}
 (a,s)=1,\quad (as,2)=1,\quad (d,s)=1,\quad (u,a)=1.
\label{eq:primitive}
\end{equation}
They imply \(a,s\) are odd.  Put
\[
 q=as,\qquad n(z)=sV(z)=aD(z)+2=su+qz.
\]
TPC-127 proves
\[
 \mu(D(z))\mu(V(z))
 =
 \lambda(q)\lambda(n(z)-2)\lambda(n(z))
 \mu^2(D(z))\mu^2(V(z))
\]
and transports the literal phase and weight without changing the
prefix order \citep{WangTPC127}.

The purpose here is narrower: replace the two squarefree indicators
by a finite family of exact congruence fibers and put every
truncation cost in a visible ledger.

\section{Exact square-divisor expansion}

For \(R\ge1\), define
\begin{equation}
 Q_R(m)=\sum_{\substack{k\le R\\k^2\mid m}}\mu(k).
\label{eq:QR}
\end{equation}
The full identity
\begin{equation}
 \mu^2(m)=\sum_{k^2\mid m}\mu(k)
\label{eq:squarefree-identity}
\end{equation}
follows prime by prime.

\begin{lemma}[Compatibility of two square divisors]
\label{lem:compatibility}
For positive integers \(k,\ell\), the simultaneous conditions
\[
 k^2\mid D(z),\qquad \ell^2\mid V(z)
\]
have an integer solution \(z\) if and only if
\begin{equation}
 (k,s)=1,\qquad(\ell,a)=1,\qquad(k,\ell)=1.
\label{eq:compatibility}
\end{equation}
When soluble, they select one residue
\[
 z\equiv r_{k,\ell}\pmod{k^2\ell^2}.
\]
\end{lemma}

\begin{proof}
If a prime \(p\) divides both \(k\) and \(s\), then
\(p\mid D(z)\) and \(p\mid s\), hence \(p\mid d\), contradicting
\((d,s)=1\).  The argument for \((\ell,a)\) is identical.  If
\(p\mid(k,\ell)\), then \(p^2\) divides both \(D(z)\) and \(V(z)\),
and consequently divides
\[
 sV(z)-aD(z)=2,
\]
which is impossible.

Conversely, if \eqref{eq:compatibility} holds, the first congruence
selects one residue modulo \(k^2\), and the second selects one
residue modulo \(\ell^2\).  Since \((k,\ell)=1\), the Chinese
remainder theorem combines them uniquely modulo \(k^2\ell^2\).
\end{proof}

\begin{theorem}[Exact truncated CRT expansion]
\label{thm:crt-expansion}
For every \(z\in I\),
\begin{equation}
\boxed{
 Q_R(D(z))Q_R(V(z))
 =
 \sum_{\substack{k,\ell\le R\\
 (k,s)=(\ell,a)=(k,\ell)=1}}
 \mu(k)\mu(\ell)
 \one_{z\equiv r_{k,\ell}\ ({\rm mod}\ k^2\ell^2)}.}
\label{eq:crt-z}
\end{equation}
Equivalently, in the \(n\)-frame each nonempty term is the unique
residue class solving
\begin{equation}
 n\equiv2\pmod{ak^2},\qquad
 n\equiv0\pmod{s\ell^2},
\label{eq:crt-n}
\end{equation}
with modulus
\begin{equation}
 q_{k,\ell}=ask^2\ell^2\le asR^4.
\label{eq:modulus-inflation}
\end{equation}
\end{theorem}

\begin{proof}
Expand the product in \eqref{eq:QR} and use
\cref{lem:compatibility}, proving \eqref{eq:crt-z}.  Since
\(aD=n-2\) and \(sV=n\), the two square divisibilities are exactly
\eqref{eq:crt-n}.  The compatibility conditions and
\((a,s)=1\) make the two displayed moduli coprime, so their product
is \eqref{eq:modulus-inflation}.
\end{proof}

\begin{remark}[No density substitution]
The coefficient \(\mu(k)\mu(\ell)\), the compatibility conditions,
and the actual CRT residue all remain in
\eqref{eq:crt-z}.  Replacing them by a squarefree density would not
be an identity and would erase the sign needed for reassembly.
\end{remark}

\section{A primitive truncation tail}

Let
\[
 Y\ge
 \max_{z\in I}\{D(z),V(z),2\},\qquad
 \tau_*(Y)=\max_{1\le m\le Y}\tau(m).
\]

\begin{lemma}[One-form square-divisor tail]
\label{lem:one-tail}
For a primitive positive affine form \(L(z)=b+cz\) on \(I\), with
\((b,c)=1\) and \(L(z)\le Y\),
\begin{equation}
 \sum_{z\in I}
 \left|\mu^2(L(z))-Q_R(L(z))\right|
 \le
 \frac NR+\sqrt Y+1.
\label{eq:one-tail}
\end{equation}
\end{lemma}

\begin{proof}
By \eqref{eq:squarefree-identity},
\[
 |\mu^2(L(z))-Q_R(L(z))|
 \le
 \sum_{\substack{k>R\\k^2\mid L(z)}}1.
\]
If \((k,c)>1\), the congruence is insoluble by primitivity.  Otherwise
it selects at most \(N/k^2+1\) integers of \(I\).  Since
\(k\le\sqrt Y\),
\[
 \sum_{R<k\le\sqrt Y}\left(\frac N{k^2}+1\right)
 \le \frac NR+\sqrt Y+1.
\]
\end{proof}

\begin{theorem}[Two-mask truncation bound]
\label{thm:tail}
For every complex sequence \(w_z\) with
\(\|w\|_\infty=\max_{z\in I}|w_z|\),
\begin{align}
 &\left|
 \sum_{z\in I}w_z
 \left\{
 \mu^2(D(z))\mu^2(V(z))
 -
 Q_R(D(z))Q_R(V(z))
 \right\}
 \right|
\notag\\
 &\qquad\le
 \|w\|_\infty(1+\tau_*(Y))
 \left(\frac NR+\sqrt Y+1\right).
\label{eq:two-tail}
\end{align}
In particular the right side is
\[
 \|w\|_\infty
 Y^{o(1)}\left(\frac NR+\sqrt Y\right).
\]
\end{theorem}

\begin{proof}
Write \(A=\mu^2(D)\), \(B=\mu^2(V)\),
\(A_R=Q_R(D)\), and \(B_R=Q_R(V)\).  Then
\[
 AB-A_RB_R=(A-A_R)B+A_R(B-B_R).
\]
Here \(|B|\le1\), while
\[
 |A_R|\le\sum_{k^2\mid D}1\le\tau(D)\le\tau_*(Y).
\]
Apply \cref{lem:one-tail} to \(D\) and \(V\), multiply by
\(\|w\|_\infty\), and use the triangle inequality.
\end{proof}

\begin{remark}[The \(\sqrt Y\) floor is real for this proof]
The \(+1\) in the count of each square-divisor progression sums to
\(O(\sqrt Y)\).  Removing that term uniformly in moving slopes and
origins would require additional information.  This paper does not
hide it in \(Y^{o(1)}N/R\).
\end{remark}

\section{Component census and exponent ledger}

\begin{proposition}[Mass-aware CRT census]
\label{prop:census}
Let \(\mathcal C_R\) be the sum, over compatible \(k,\ell\le R\), of
the number of \(z\in I\) in the corresponding CRT fiber.  Then
\begin{equation}
 \mathcal C_R
 \le
 N\left(\sum_{k\ge1}\frac1{k^2}\right)^2+R^2
 \ll N+R^2.
\label{eq:census}
\end{equation}
\end{proposition}

\begin{proof}
One residue modulo \(k^2\ell^2\) contributes at most
\(N/(k^2\ell^2)+1\) points.  Sum this upper bound and discard the
compatibility restrictions.
\end{proof}

The census matters because a uniform estimate relative to each
component's \emph{own} absolute mass costs \(N+R^2\), not the
wasteful \(R^2N\).  Physical weights still have to be included in
the corresponding weighted census; \cref{prop:census} alone does
not prove that outer bound.

Suppose, for an exponent audit, that
\[
 N=X^{\nu+o(1)},\quad
 Y\le X^{\upsilon+o(1)},\quad
 as\le X^{\kappa+o(1)},\quad
 R=X^{\rho+o(1)},
\]
and put
\begin{equation}
 R_{\rm eff}=\min\{R,\sqrt Y\},
 \qquad
 \rho_{\rm eff}=\min\{\rho,\upsilon/2\}.
\label{eq:effective-cutoff}
\end{equation}
Only \(k,\ell\le R_{\rm eff}\) can occur on the interval.  Every retained
correlation theorem must be uniform up to modulus exponent
\begin{equation}
 \kappa+4\rho_{\rm eff}.
\label{eq:uniformity-cost}
\end{equation}

To state a relative tail exponent, the comparison mass must be
declared.  Let \(\mathcal V_X\) be the actual absolute mass of the
untruncated literal block and assume
\[
 \|w\|_\infty=X^{o(1)},
 \qquad
 \mathcal V_X=X^{\omega+o(1)}.
\]
If \(R<\sqrt Y\), \cref{thm:tail} gives
\begin{equation}
 \eta_{\rm tail}
 =
 \min\left\{\omega-\nu+\rho,\,
             \omega-\frac{\upsilon}{2}\right\},
\label{eq:tail-exponent}
\end{equation}
provided the displayed minimum is positive.  In the common
mass-comparable case \(\omega=\nu\), this reduces to
\(\min\{\rho,\nu-\upsilon/2\}\).  If \(R\ge\sqrt Y\), then
\(Q_R(m)=\mu^2(m)\) for every \(m\le Y\), so the truncation tail is
identically zero rather than being governed by
\eqref{eq:tail-exponent}.

For each retained CRT component \(c\), let \(\mathcal V_c\) be its
actual absolute mass.  Define the \emph{weighted census cost} by the
theorem-backed inequality
\begin{equation}
 \sum_c\mathcal V_c
 \le X^{\ell_{\rm cen}+o(1)}\mathcal V_X.
\label{eq:weighted-census}
\end{equation}
If the literal weights are uniformly bounded and
\(\mathcal V_X=X^{\nu+o(1)}\), the point census gives the crude choice
\begin{equation}
 \ell_{\rm cen}=(2\rho_{\rm eff}-\nu)_+.
\label{eq:census-cost}
\end{equation}
Without those mass hypotheses, \eqref{eq:census-cost} is unavailable
and \(\ell_{\rm cen}\) must come from the actual weighted archive.

If each CRT component saves \(X^{-\eta_{\rm corr}}\) relative to
its own \(\mathcal V_c\), the certificate returns
\begin{equation}
 \boxed{
 \eta_{\rm block}^{\rm cert}
 =
 \begin{cases}
 \min\{\eta_{\rm tail},
       \eta_{\rm corr}-\ell_{\rm cen}\},
       &R<\sqrt Y,\\
 \eta_{\rm corr}-\ell_{\rm cen},
       &R\ge\sqrt Y.
 \end{cases}}
\label{eq:block-certificate}
\end{equation}
The formula is relative to the declared actual absolute mass.  It is
a loss ledger, not an assertion of the correlation or weighted-census
inputs.  It is a positive saving only when
\(\eta_{\rm tail}>0\) and
\(\eta_{\rm corr}>\ell_{\rm cen}\) in the first regime, or when
\(\eta_{\rm corr}>\ell_{\rm cen}\) in the exact-cutoff regime.

\section{Stop/go rules and claim boundary}

\paragraph{TAIL STOP.}
If \(\nu\le\upsilon/2\), the elementary uniform tail
\eqref{eq:two-tail} gives no positive relative power in the
mass-comparable regime while \(R<\sqrt Y\).  One must either prove a
stronger weighted tail theorem, choose the exact cutoff
\(R\ge\sqrt Y\), or use a different decomposition.  The exact cutoff
removes the tail but caps the effective component range at
\(R_{\rm eff}=\sqrt Y\), with modulus and census costs
\(\kappa+2\upsilon\) and potentially
\((\upsilon-\nu)_+\), respectively.

\paragraph{MODULUS STOP.}
Choosing \(R\) large enough to improve the tail inflates the actual
modulus to \(asR^4\).  A theorem uniform only for fixed modulus is
not callable on these fibers.

\paragraph{CENSUS STOP.}
The bound \eqref{eq:census} is unweighted.  If the literal physical
weights, the active absolute mass, or outer keys do not satisfy the
mass-comparable hypotheses, the actual bound
\eqref{eq:weighted-census} must replace
\eqref{eq:census-cost}.

\paragraph{ARITHMETIC GO.}
Proceed only with a proved estimate for the actual shift-two
Liouville correlation on every retained CRT fiber, with the original
phase, weight, origin, and prefix quantifiers.  Such a theorem would
be genuine \(\Ltwo\) input.

\begin{center}
\small
\begin{tabularx}{\textwidth}{@{}p{0.47\textwidth}X@{}}
\toprule
Statement & Level and status\\
\midrule
Compatibility lemma and CRT expansion
& Exact finite theorem (\(\Lzero\)).\\
Primitive tail \eqref{eq:two-tail} and census
\eqref{eq:census}
& Unconditional elementary bounds (\(\Lzero\)).\\
Attachment to the TPC-127 actual fixed-\(2\) frame
& Literal interface (\(\Lone\)).\\
Weighted outer census on the complete family
& Not proved here.\\
Uniform cancellation on retained CRT fibers
& Not proved; arithmetic gate (\(\Ltwo\) if proved).\\
H3, zero-mode saving, or endpoint theorem
& Not proved.\\
\bottomrule
\end{tabularx}
\end{center}

TPC-128 opens the squarefree masks exactly and quantifies the price:
tail improvement, modulus growth, and component census compete.
Squarefree density by itself supplies none of the required signed
correlation.  No parity breach, prime-pair lower bound, or twin-prime
conclusion is claimed.

\begingroup
\footnotesize
\raggedright
\bibliographystyle{plainnat}
\bibliography{references}
\endgroup

\end{document}
